\documentclass[11pt]{article}
\usepackage{explorer}
\newcommand{\ExplorerLanguage}{FR}
\newcommand{\ExplorerNumber}{11}
\newcommand{\ExplorerKind}{Solutions détaillées}
\begin{document}
\ExplorerTitle{Lycée · optimisation et olympiades}{Inégalités et cas d’égalité}
\ExplorerSection{Exercice 1}
Pour \(x>0\), prouver \(x+\frac{1}{x}\geq 2\) et déterminer l’égalité.\par
\(x+\frac{1}{x}-2=\frac{(x-1)^{2}}{x}\geq 0\). L’égalité a lieu exactement pour \(x=1\).\par
\ExplorerSection{Exercice 2}
Déterminer le minimum de \(x^{2}-6x+13\) sur \(\mathbb{R}\).\par
\(x^{2}-6x+13=(x-3)^{2}+4\geq 4\). Le minimum 4 est atteint en \(x=3\).\par
\ExplorerSection{Exercice 3}
Un rectangle a pour périmètre 24 m. Quelle est son aire maximale ?\par
Ses côtés \(a,b>0\) vérifient \(a+b=12\). Donc \(ab\leq 36\), avec égalité pour \(a=b=6\,\mathrm{m}\). L’aire maximale est \(36\,\mathrm{m}^{2}\).\par
\ExplorerSection{Exercice 4}
Pour a,b réels, prouver \(a^{2}+b^{2}\geq 2ab\). Faut-il supposer a et b positifs ?\par
\(a^{2}+b^{2}-2ab=(a-b)^{2}\geq 0\) pour tous réels a,b. Aucune positivité n’est nécessaire ; l’égalité équivaut à \(a=b\).\par
\ExplorerSection{Problème : raisonnement complet}
On dispose de 40 m de clôture pour trois côtés d’un enclos rectangulaire adossé à un mur. La largeur x, perpendiculaire au mur, est strictement positive. Exprimer la longueur et l’aire en fonction de x, déterminer le domaine, puis prouver l’aire maximale sans dérivée. Comparer avec un rectangle fermé de périmètre 40 m.\par
La longueur vaut \(40-2x\) et \(0<x<20\). L’aire est \(A(x)=x(40-2x)=200-2(x-10)^{2}\leq 200\). L’égalité est atteinte pour largeur 10 m et longueur 20 m : \(200\,\mathrm{m}^{2}\). Pour quatre côtés clôturés, \(a+b=20\) et l’aire maximale est \(100\,\mathrm{m}^{2}\), atteinte par le carré de côté 10 m.\par
\end{document}
